\documentclass[a4paper, serif, 10pt]{moderncv}

%% moderncv
% color and style
\moderncvcolor{black}
\moderncvstyle{banking}

%% page layout/margins
\usepackage[top=2.5cm, bottom=2.5cm, left=1.5cm, right=1.5cm]{geometry}

\nopagenumbers{}

%% Babel config for Dutch language
% ref:
% - https://latex3.github.io/babel/guides/locale-dutch.html
\usepackage[dutch]{babel}

%% fontspec
\usepackage{fontspec}

\IfFontExistsTF{Linux Libertine}{
  \setmainfont[Ligatures={TeX, Common}, Numbers=OldStyle]{Linux Libertine}
}{}
\IfFontExistsTF{Linux Libertine O}{
  \setmainfont[Ligatures={TeX, Common}, Numbers=OldStyle]{Linux Libertine O}
}{}

%% CTeX
% CJK support, used to show CJK characters in the resume
%
% - heading=false: disable ctex heading
% - punct=kaiming: use kaiming punctuations styles for CJK
% - scheme=plain: use plain scheme, do not override \normalsize font size
% - space=auto: space settings for CJK characters
%
% CJK fonts are resolved per locale language, e.g. Noto Serif/Sans CJK JP
% for Japanese resumes, with fallback to the next region if not installed
%
% ref:
% - http://ctan.mirrorcatalogs.com/language/chinese/ctex/ctex.pdf
\usepackage[UTF8, heading=false, punct=kaiming, scheme=plain, space=auto]{ctex}

\IfFontExistsTF{Noto Serif CJK SC}{
  \setCJKmainfont{Noto Serif CJK SC}
}{}
\IfFontExistsTF{Noto Sans CJK SC}{
  \setCJKsansfont{Noto Sans CJK SC}
}{}

%% line spacing
\usepackage{setspace}
\setstretch{1.125}

%% URL styling - use normal text instead of monospace
\urlstyle{same}

%% Patch \httplink and \httpslink to handle full URLs
% The original moderncv macros always prepend http:// or https:// to URLs.
% This causes issues when the URL already has a protocol (e.g., https://example.com)
% resulting in malformed URLs like https://https://example.com.
% This patch checks if the URL already contains "://" and uses it directly if so.
% Note: We use \str_set:Nx (x-expansion) to expand macros like \@homepage first.
\ExplSyntaxOn
\str_new:N \g_yamlresume_url_str

\RenewDocumentCommand{\httpslink}{O{}m}{
  \str_set:Nx \g_yamlresume_url_str {#2}
  \str_if_in:NnTF \g_yamlresume_url_str {://}
    { \url{#2} }
    { \url{https://#2} }
}

\RenewDocumentCommand{\httplink}{O{}m}{
  \str_set:Nx \g_yamlresume_url_str {#2}
  \str_if_in:NnTF \g_yamlresume_url_str {://}
    { \url{#2} }
    { \url{http://#2} }
}
\ExplSyntaxOff

\name{Jeroen van Leeuwen}{}
\title{Business Analist met focus op datagedreven procesverbetering}
\phone[mobile]{+31 6 12345678}
\email{jeroen.vanleeuwen@analytiq.nl}
\homepage{https://jeroenvanleeuwen.nl}

\address{Keizersgracht 123, Amsterdam, Noord-Holland, Nederland, 1015 CJ}{}{}

\extrainfo{{\small \faLinkedin}\ \href{https://www.linkedin.com/in/jeroenvanleeuwen}{@jeroenvanleeuwen} {} {} {} • {} {} {} 
{\small \faMedium}\ \href{https://medium.com/@jeroenvanleeuwen}{@jeroenvanleeuwen}}

\begin{document}

\maketitle

\section{Basis Informatie}

\cvline{}{\begin{itemize}
\item Business analist met 7+ jaar ervaring in de financiële dienstverlening, retail en publieke sector
\item Sterk in het vertalen van bedrijfsdoelen naar heldere requirements, procesmodellen en data-oplossingen
\item Ervaren in agile teams, stakeholder management en het begeleiden van multidisciplinaire workshops
\item Combineert analytisch vermogen met communicatieve vaardigheden om draagvlak te creëren voor verandering
\end{itemize}}
\section{Opleidingen}

\cventry{sep 2014–aug 2016}
        {Master, Bedrijfskunde, Score: 8.2}
        {Universiteit van Amsterdam}
        {\url{https://www.uva.nl}}
        {}
        {\begin{itemize}
\item Afgestudeerd met specialisatie in digitaal transformatiemanagement
\item Scriptie over datagedreven procesverbetering in de financiële sector
\item Actief lid van de studievereniging en organisator van carrière-evenementen
\end{itemize}
\textbf{Cursussen}: Strategisch Management, Business Process Management, Data-analyse voor Bedrijfskunde, Informatiesystemen, Financieel Management, Operationeel Onderzoek}

\cventry{sep 2011–jun 2014}
        {Bachelor, Informatiekunde, Score: 7.8}
        {Universiteit Utrecht}
        {\url{https://www.uu.nl}}
        {}
        {\begin{itemize}
\item Basis gelegd in informatiesystemen, data modellering en organisatieprocessen
\item Minor gevolgd in Bedrijfseconomie aan de Universiteit van Amsterdam
\item Stage gelopen bij een gemeentelijke afdeling voor digitale dienstverlening
\end{itemize}
\textbf{Cursussen}: Informatiesystemen, Databases, Systeemontwikkeling, Organisatiekunde, Statistiek, Projectmanagement}
\section{Werkervaring}

\cventry{jan 2022–Heden}
        {Senior Business Analist}
        {NN Group}
        {\url{https://www.nn-group.com}}
        {}
        {\begin{itemize}
\item Verantwoordelijk voor requirements engineering en impactanalyse voor digitale claimprocessen
\item Samenwerken met product owners, architecten en compliance om wet- en regelgeving te vertalen naar oplossingen
\item Opstellen van user stories, acceptatiecriteria en procesmodellen (BPMN) voor agile teams
\item Sturen op datakwaliteit en het definiëren van KPI's voor continue procesverbetering
\item Begeleiden van twee junior business analisten en verzorgen van interne trainingen
\end{itemize}
\textbf{Trefwoorden}: Requirements Engineering, Procesoptimalisatie, Data-analyse, Agile, Verzekeringen}

\cventry{mrt 2018–dec 2021}
        {Business Analist}
        {Coolblue}
        {\url{https://www.coolblue.nl}}
        {}
        {\begin{itemize}
\item Analyseerde klant- en orderprocessen en vertaalde knelpunten naar concrete verbetervoorstellen
\item Werkte nauw samen met e-commerce, logistiek en klantenservice om end-to-end processen te optimaliseren
\item Stelde dashboards op in Power BI voor het monitoren van operationele prestaties
\item Faciliteerde workshops met stakeholders om requirements en prioriteiten te bepalen
\item Droeg bij aan de invoering van een nieuwe order management systeem
\end{itemize}
\textbf{Trefwoorden}: eCommerce, Procesanalyse, Power BI, Stakeholdermanagement, Scrum}

\cventry{sep 2016–feb 2018}
        {Junior Business Analist}
        {Capgemini}
        {\url{https://www.capgemini.com}}
        {}
        {\begin{itemize}
\item Ondersteunde senior analisten bij het in kaart brengen van bedrijfsprocessen en informatiesystemen
\item Maakte procesbeschrijvingen, datastroomdiagrammen en testscenario's
\item Participeerde in klantworkshops en vertaalde bevindingen naar heldere rapportages
\item Leerde werken volgens de Waterval- en Agile-methodiek binnen overheidsprojecten
\end{itemize}
\textbf{Trefwoorden}: Business Analyse, Procesmodellering, Testscenario's, Overheid, Agile}
\section{Talen}

\cvline{Nederlands}{Moedertaal of tweetalige vaardigheid \hfill \textbf{Trefwoorden}: Moedertaal}
\cvline{Engels}{Volledige professionele vaardigheid \hfill \textbf{Trefwoorden}: TOEFL iBT 108, Zakelijk Engels}
\cvline{Duits}{Beperkte werkvaardigheid \hfill \textbf{Trefwoorden}: Goethe B1}
\section{Vaardigheden}

\cvline{Business Analyse}{Expert \hfill \textbf{Trefwoorden}: Requirements Engineering, Procesmodellering, Stakeholdermanagement, User Stories, BPMN}
\cvline{Data \& Analytics}{Geavanceerd \hfill \textbf{Trefwoorden}: SQL, Power BI, Tableau, Excel, Python}
\cvline{Agile \& Scrum}{Geavanceerd \hfill \textbf{Trefwoorden}: Scrum, Kanban, SAFe, Jira, Confluence}
\cvline{Domeinkennis}{Geavanceerd \hfill \textbf{Trefwoorden}: Verzekeringen, Retail, Financiële dienstverlening, eCommerce}
\cvline{Tools}{Intermediair \hfill \textbf{Trefwoorden}: Visio, Azure DevOps, Miro, Figma, SAP}
\section{Onderscheidingen}

\cventry{jun 2016}
        {Universiteit van Amsterdam}
        {Beste Afstudeerscriptie Bedrijfskunde}
        {}
        {}
        {Toegekend voor een scriptie over datagedreven procesverbetering in de financiële sector, met een beoordeling van 9,0.}

\cventry{nov 2020}
        {Coolblue}
        {Coolblue Innovation Award}
        {}
        {}
        {Erkend voor het ontwikkelen van een nieuw KPI-dashboard dat de orderdoorlooptijd met 15\% verkortte.}
\section{Certificaten}

\cventry{mrt 2021}
        {Scrum.org}
        {Professional Scrum Master I}
        {\url{https://www.scrum.org/professional-scrum-certifications/professional-scrum-master-i-assessment}}
        {}
        {}

\cventry{sep 2022}
        {IIBA}
        {CBAP}
        {\url{https://www.iiba.org/certifications/certified-business-analysis-professional/}}
        {}
        {}

\cventry{feb 2019}
        {PeopleCert}
        {PRINCE2 Foundation}
        {\url{https://www.peoplecert.org}}
        {}
        {}
\section{Publicaties}

\cventry{nov 2021}
        {Datagedreven procesverbetering in de verzekeringssector}
        {Tijdschrift voor Business Analyse}
        {\url{https://www.tijdschriftbusinessanalyse.nl/artikel/datagedreven-procesverbetering}}
        {}
        {\begin{itemize}
\item Beschrijft een praktijkgerichte aanpak voor het combineren van procesanalyse en data-analyse
\item Gaat in op KPI-gedreven sturing en stakeholdermanagement in verzekeringsprocessen
\item Inclusief checklist voor business analisten die datagedreven verbeteringen willen doorvoeren
\end{itemize}}
\section{Projecten}

\cventry{jan 2023–okt 2023}
        {Herontwerp van de klantreis voor schadeclaims van eerste melding tot uitbetaling.}
        {Klantreisoptimalisatie NN Group}
        {\url{https://www.nn-group.com/innovatie/klantreis}}
        {}
        {\begin{itemize}
\item Voerde klantinterviews en journey mapping-workshops uit met vertegenwoordigers van klantenservice, schade en IT
\item Identificeerde vijf kritieke knelpunten en vertaalde deze naar een prioriteitenlijst met quick wins
\item Realiseerde een verkorting van de gemiddelde doorlooptijd met 20\% door processtappen te elimineren
\end{itemize}
\textbf{Trefwoorden}: Klantreis, Procesoptimalisatie, Design Thinking, Data-analyse}

\cventry{feb 2020–nov 2021}
        {Migratie van een verouderd datawarehouse naar een cloudgebaseerde analytics-omgeving.}
        {Datawarehouse migratie Coolblue}
        {\url{https://www.coolblue.nl/tech/datawarehouse}}
        {}
        {\begin{itemize}
\item Vertaalde rapportagebehoeften van finance, logistiek en marketing naar datamodellen
\item Schreef functionele specificaties en testcases voor ETL-processen
\item Begeleidde gebruikers bij de acceptatie en adoptie van nieuwe Power BI-rapporten
\end{itemize}
\textbf{Trefwoorden}: Datawarehouse, ETL, Power BI, Requirements}
\section{Interesses}

\cvline{Hardlopen}{Halve marathon, Trailrunning}
\cvline{Culinair}{Koken, Wijnproeven}
\cvline{Technologie}{AI, Python, Open data}
\section{Vrijwilligerswerk}

\cventry{sep 2020–Heden}
        {Mentor}
        {Stichting Business Analyse Nederland}
        {\url{https://www.businessanalyse.nl}}
        {}
        {\begin{itemize}
\item Begeleidt startende business analisten bij hun professionele ontwikkeling
\item Geeft workshops over requirements engineering en stakeholdercommunicatie
\item Draagt bij aan de jaarlijkse kennisdag voor de beroepsvereniging
\end{itemize}}
    

\cventry{jan 2019–dec 2021}
        {Vrijwilliger Data Coach}
        {Amsterdam Data Collective}
        {\url{https://www.amsterdamdatacollective.nl}}
        {}
        {\begin{itemize}
\item Ondersteunde non-profitorganisaties bij het opzetten van eenvoudige data-analyses
\item Hielp bij het vertalen van maatschappelijke vraagstukken naar meetbare indicatoren
\item Verzorgde trainingen in Excel en Power BI voor vrijwilligers
\end{itemize}}
    
\end{document}